\documentclass[12pt,a4paper]{article}

\usepackage{amsmath, amssymb, amsthm}
\usepackage{geometry}
\geometry{margin=2.5cm}

\title{Modelarea Matematică a Sistemelor Materiale \\ \large Analiza și Simularea Modelelor}
\author{}
\date{}

\begin{document}

\maketitle

\section{Analiza și simularea modelelor}

Analiza și simularea modelelor matematice reprezintă etape esențiale în înțelegerea comportamentului sistemelor materiale. Ele permit investigarea soluțiilor, stabilității, regimurilor de funcționare și sensibilității la variația parametrilor.

\subsection{6.1. Analiză analitică}

Analiza analitică urmărește obținerea de soluții exacte sau proprietăți fundamentale ale modelului.

\subsubsection*{Soluții exacte}
Pentru unele modele simple (liniare, cu coeficienți constanți), soluțiile pot fi obținute explicit. De exemplu:


\[
\frac{dx}{dt} = ax \quad \Rightarrow \quad x(t) = x_0 e^{at}.
\]



\subsubsection*{Stabilitate}
Stabilitatea soluțiilor este analizată prin:
\begin{itemize}
    \item studiul valorilor proprii ale sistemelor liniare;
    \item analiza Lyapunov pentru sisteme neliniare;
    \item determinarea punctelor de echilibru.
\end{itemize}

\subsubsection*{Regim staționar vs. tranzitoriu}
Un model poate prezenta:
\begin{itemize}
    \item \textbf{regim staționar}: soluții independente de timp;
    \item \textbf{regim tranzitoriu}: evoluție temporară până la stabilizare.
\end{itemize}

\subsection{6.2. Analiză numerică}

Pentru majoritatea modelelor reale, soluțiile analitice nu sunt disponibile, fiind necesare metode numerice.

\subsubsection*{Metode Euler și Runge--Kutta}
Pentru ODE:


\[
x_{n+1} = x_n + h f(x_n, t_n) \quad \text{(Euler)}.
\]


Metodele Runge--Kutta de ordin superior oferă precizie mai mare.

\subsubsection*{Diferențe finite}
Aproximarea derivatelor:


\[
\frac{\partial u}{\partial x} \approx \frac{u_{i+1} - u_i}{\Delta x}.
\]



\subsubsection*{Elemente finite}
Metodă puternică pentru PDE complexe, bazată pe:
\begin{itemize}
    \item împărțirea domeniului în elemente;
    \item aproximarea soluției prin funcții de bază;
    \item formularea slabă a problemei.
\end{itemize}

\subsubsection*{Discretizare spațială și temporală}
Modelul continuu este transformat într-un sistem discret:


\[
u(x,t) \rightarrow u_{i}^{n}.
\]



\subsection{6.3. Validare și calibrare}

Validarea verifică dacă modelul reproduce realitatea.

\subsubsection*{Ajustarea parametrilor}
Parametrii sunt calibrați pentru a minimiza diferența dintre model și date.

\subsubsection*{Compararea cu măsurători}
Se utilizează:
\begin{itemize}
    \item erori absolute și relative;
    \item norme $L^2$;
    \item analize statistice.
\end{itemize}

\subsubsection*{Analiza erorilor}
Evaluarea:
\begin{itemize}
    \item erorilor de modelare;
    \item erorilor numerice;
    \item incertitudinii în date.
\end{itemize}

\end{document}
